\documentclass[12pt]{article}
\usepackage[utf8]{inputenc}
\usepackage[T1]{fontenc}
\usepackage{hyperref}
\title{Uma Language Resources for K-12 Curriculum Development}
\author{Compiled by Claude Sonnet 5.5 (claude-sonnet-5-5) for Andr\'as Kornai}
\date{}
\begin{document}
\maketitle

\section{Introduction}
Uma (native name Pipikoro; ISO 639-3 ppk) is an Austronesian language of Central and South Sulawesi, Indonesia. Search results turned up scripture and audio materials and general reference pages, but no school textbooks.

\section{Primary Educational Resources}
No primer or textbook was found in search results.

\section{Linguistic and Reference Materials}
Search summaries mention Michael Martens' ``Notes on Uma Verbs'' and ``Uma Dialect Word Lists'' (2014); these were not located with direct URLs and are not cited as links.
\begin{itemize}
\item Wikipedia, Uma language: \url{https://en.wikipedia.org/wiki/Uma_language}
\item Omniglot, Uma: \url{https://omniglot.com/writing/uma.htm}
\item Endangered Languages Project sample: \url{https://endangeredlanguages.com/lang/1361/samples/11600}
\end{itemize}

\section{Religious and Community Resources}
\begin{itemize}
\item ScriptureEarth, Uma (ppk) scripture resources, including a Bible text and a Gospel of Mark film: \url{https://www.scriptureearth.org/00eng.php?iso=ppk}
\end{itemize}

\section{Digital and Audio Resources}
\begin{itemize}
\item Global Recordings Network, Uma audio: \url{https://globalrecordings.net/en/language/ppk}
\end{itemize}

\section{Curriculum Development Context}
Uma is still spoken by all adults but not all children, and is taught as a subject in some schools. The scripture texts give a written standard that could seed reading materials, but no dedicated educational materials were found.
\end{document}